%% FullCourse_10Lecs -- Lecture 9, Topic 9.4: Project Organization
%% Source: ./archive_pre_split/Lec09_Agentic_AI_v3.tex (section 8)
%% Pass B: layer in refinements from
%%         ../../MiniCourse_8hr/Slides/Module4_Agentic_AI/M4_T7_Specs_Harness_Skills.tex
%%         ../../MiniCourse_8hr/Slides/Module4_Agentic_AI/M4_T8_Remote_HPC_Subagents.tex
%% Rewritten 2026-07-06: Act 4 "Management" of the five-act structure. Opens with the
%%         researcher-as-project-manager paradigm shift; HA-Ramsey milestone ladder,
%%         division-of-labor contract, quality gates, orchestration (Lec10 T8 F19/F20/
%%         F22/F25/F38); fellows provenance rules as rules+hooks example (Lec10 T4).
%%         HPC segment compressed 4->2 frames. See Lec09_Rewrite_Plan_2026-07-06.md.

%% Shared preamble for the FullCourse_10Lecs topic decks.
%%
%% Each topic .tex \input's this file, then defines its own \title, \date
%% override (if any), \graphicspath, and \begin{document}.
%%
%% Reconstructed verbatim from the inlined copy preserved in
%% Lec10_Case_Studies/Lec10_T8_Case6_DeepRamsey.tex (lines 23-190), which is
%% explicitly annotated there as "from ../shared_preamble.tex, copied verbatim".
%% Base preamble matches MiniCourse_8hr/Slides/shared_preamble.tex; this version
%% additionally provides \sectiondivider, the conceptbox/cautionbox/examplebox/
%% takeawaybox tcolorboxes, and \compacttables used across the split decks.

\documentclass[10pt,english,aspectratio=169]{beamer}

\usepackage{ae,aecompl}
\usepackage[T1]{fontenc}
\usepackage{textcomp}
\usepackage[utf8]{inputenc}
\usepackage{babel}
\usepackage{datetime2}

\setcounter{secnumdepth}{3}
\setcounter{tocdepth}{3}

\usepackage{amsmath, amssymb, amsthm, graphicx}
\usepackage{booktabs, multirow}
\usepackage{tikz}
\usepackage{pgfplots}
\pgfplotsset{compat=1.16}
\usepackage[most]{tcolorbox}
\tcbuselibrary{skins, breakable, theorems}
\usepackage{listings}
\usepackage{xcolor}

\lstset{
  basicstyle=\ttfamily\small,
  keywordstyle=\color{blue!80!black}\bfseries,
  commentstyle=\color{green!50!black}\itshape,
  stringstyle=\color{orange!80!black},
  backgroundcolor=\color{gray!8},
  frame=single,
  rulecolor=\color{gray!40},
  breaklines=true,
  showstringspaces=false,
  numbers=none,
}

\lstdefinestyle{pythonstyle}{
    language=Python,
    basicstyle=\ttfamily\footnotesize,
    keywordstyle=\color{blue!80!black}\bfseries,
    commentstyle=\color{green!50!black}\itshape,
    stringstyle=\color{orange!80!black},
    frame=single,
    breaklines=true,
    showstringspaces=false,
    numbers=left,
    numbersep=5pt,
    numberstyle=\tiny\color{gray},
    backgroundcolor=\color{gray!8},
    rulecolor=\color{gray!40}
}

\ifx\hypersetup\undefined
  \AtBeginDocument{%
    \hypersetup{bookmarks=false,breaklinks=true,pdfborder={0 0 0},colorlinks=true,
      linkcolor=DarkText,urlcolor=RedTitle}%
  }
\else
  \hypersetup{bookmarks=false,breaklinks=true,pdfborder={0 0 0},colorlinks=true,
    linkcolor=DarkText,urlcolor=RedTitle}%
\fi

\makeatletter
\newcommand\makebeamertitle{%
  \begin{frame}[plain]
    \vspace*{1.0cm}
    \centering
    {\usebeamerfont{title}\usebeamercolor[fg]{title}\inserttitle\par}
    \vspace{1.0cm}
    {\usebeamerfont{author}\usebeamercolor[fg]{author}\insertauthor\par}
    \vspace{0.2cm}
    {\usebeamerfont{date}\usebeamercolor[fg]{date}\insertdate\par}
  \end{frame}}%
\AtBeginDocument{%
  \let\origtableofcontents=\tableofcontents
  \def\tableofcontents{\@ifnextchar[{\origtableofcontents}{\gobbletableofcontents}}
  \def\gobbletableofcontents#1{\origtableofcontents}%
}

\usetheme{metropolis}
\usefonttheme{professionalfonts}

\definecolor{LightGray}{RGB}{230,230,230}
\definecolor{RedTitle}{RGB}{0,0,200}
\definecolor{VeryLightGray}{RGB}{245,245,245}
\definecolor{DarkText}{RGB}{0,0,0}

\setbeamercolor{background canvas}{bg=white}
\setbeamercolor{title}{fg=RedTitle}
\setbeamercolor{author}{fg=DarkText}
\setbeamercolor{date}{fg=DarkText}
\setbeamercolor{frametitle}{bg=VeryLightGray,fg=DarkText}
\setbeamercolor{normal text}{fg=DarkText}
\setbeamerfont{title}{size=\large}

\setbeamersize{text margin left=5mm, text margin right=5mm}
\addtobeamertemplate{frametitle}{}{\vspace{-0.5em}}
\newcommand{\reducevspace}{\vspace{-0.5cm}}

% Shared section divider and box styles for split topic decks.
\newcommand{\sectiondivider}[2]{%
\begin{frame}[plain,noframenumbering]
\begin{center}
\vfill
{\Large\textbf{#1}\par}
\vspace{0.35cm}
{\normalsize #2\par}
\vfill
\end{center}
\end{frame}
}

\newtcolorbox{conceptbox}[1][]{
  colback=blue!3,
  colframe=RedTitle!55,
  boxrule=0.35mm,
  arc=1mm,
  left=2mm,
  right=2mm,
  top=1mm,
  bottom=1mm,
  #1
}
\newtcolorbox{cautionbox}[1][]{
  colback=red!3,
  colframe=red!50!black,
  boxrule=0.35mm,
  arc=1mm,
  left=2mm,
  right=2mm,
  top=1mm,
  bottom=1mm,
  #1
}
\newtcolorbox{examplebox}[1][]{
  colback=green!3,
  colframe=green!45!black,
  boxrule=0.35mm,
  arc=1mm,
  left=2mm,
  right=2mm,
  top=1mm,
  bottom=1mm,
  #1
}
\newtcolorbox{takeawaybox}[1][]{
  colback=VeryLightGray,
  colframe=RedTitle!55,
  boxrule=0.35mm,
  arc=1mm,
  left=2mm,
  right=2mm,
  top=1mm,
  bottom=1mm,
  #1
}
\newcommand{\compacttables}{\renewcommand{\arraystretch}{1.15}}

\setbeamertemplate{navigation symbols}{}
\setbeamertemplate{footline}{%
  \hfill%
  \usebeamercolor[fg]{page number in head/foot}%
  \usebeamerfont{page number in head/foot}%
  \insertframenumber\,/\,\inserttotalframenumber\kern1em\vskip2pt%
}

\makeatother

\author{Zhigang Feng}
% "date with time": compile timestamp so each build is identifiable.
\date{\today\ \DTMcurrenttime}


% Suppress redundant auto-generated metropolis section page; \sectiondivider frames are the dividers we keep.
\metroset{sectionpage=none}

\graphicspath{{./pic/}{./figures/}{../../MiniCourse_8hr/Slides/Module4_Agentic_AI/pic/}{../}}


\usetikzlibrary{positioning,arrows.meta,shapes,calc}

%% Lec09 shared MAP slide: "one map, five acts" recurring diagram.
%% Input this file AFTER ../shared_preamble.tex (needs RedTitle, VeryLightGray)
%% and call \LecNineMap{k} inside a frame, where k in {1..5} highlights the
%% current act ("you are here") and k = 0 renders the closing reprise with all
%% five acts checked (used by T5's "Your Job Description" frame).
%% Filename deliberately does not match the build glob Lec*_T*.tex, so the
%% build script never tries to compile it standalone.

\newcommand{\LecNineMap}[1]{%
% Precompute per-box style selectors and the checkmark token; TikZ option
% lists are comma-split before TeX conditionals run, so \ifnum must not sit
% inside a [...] options list.
\ifnum#1=0
  \tikzset{lecmapa/.style={lecmapdone}, lecmapb/.style={lecmapdone},
           lecmapc/.style={lecmapdone}, lecmapd/.style={lecmapdone},
           lecmape/.style={lecmapdone}}%
  % green fill + the "all five acts complete" banner mark completion; a per-box
  % checkmark overflows the MANAGEMENT box, so keep the token empty here too.
  \def\lecmapchk{}%
\else
  \tikzset{lecmapa/.style={}, lecmapb/.style={}, lecmapc/.style={},
           lecmapd/.style={}, lecmape/.style={}}%
  \def\lecmapchk{}%
  \ifnum#1=1 \tikzset{lecmapa/.style={lecmapcur}}\fi
  \ifnum#1=2 \tikzset{lecmapb/.style={lecmapcur}}\fi
  \ifnum#1=3 \tikzset{lecmapc/.style={lecmapcur}}\fi
  \ifnum#1=4 \tikzset{lecmapd/.style={lecmapcur}}\fi
  \ifnum#1=5 \tikzset{lecmape/.style={lecmapcur}}\fi
\fi
\begin{center}
\begin{tikzpicture}[
    act/.style={rectangle, rounded corners, draw=black!60, fill=VeryLightGray,
        minimum width=2.6cm, minimum height=0.92cm, text width=2.5cm,
        align=center, font=\scriptsize, text=DarkText},
    lecmapcur/.style={draw=RedTitle, very thick, fill=orange!20},
    lecmapdone/.style={draw=green!45!black, thick, fill=green!10},
    q/.style={font=\tiny\itshape, text width=2.7cm, align=center, text=black!70},
    sat/.style={rectangle, rounded corners, draw=black!45, dashed, fill=white,
        text width=5.0cm, align=center, font=\tiny, text=black!70,
        minimum height=0.5cm},
    barlab/.style={font=\tiny\bfseries, text=black!60},
    arr/.style={-{Stealth[length=2mm]}, thick, gray!60}
]
    % --- five act boxes ---
    \node[act, lecmapa] (a1) at (0,0)
        {\textbf{9.1 MAP}\lecmapchk\\ \tiny what \& why};
    \node[act, lecmapb] (a2) at (2.85,0)
        {\textbf{9.2 MACHINE}\lecmapchk\\ \tiny under the hood};
    \node[act, lecmapc] (a3) at (5.7,0)
        {\textbf{9.3 METHOD}\lecmapchk\\ \tiny homotopy workflow};
    \node[act, lecmapd] (a4) at (8.55,0)
        {\textbf{9.4 MANAGEMENT}\lecmapchk\\ \tiny run the project};
    \node[act, lecmape] (a5) at (11.4,0)
        {\textbf{9.5 MINDSET}\lecmapchk\\ \tiny your role};
    \draw[arr] (a1) -- (a2);
    \draw[arr] (a2) -- (a3);
    \draw[arr] (a3) -- (a4);
    \draw[arr] (a4) -- (a5);
    % --- the student question each act answers ---
    \node[q] at (0,-0.82)    {what is this,\\ and why care?};
    \node[q] at (2.85,-0.82) {how does it work?\\ how do I start?};
    \node[q] at (5.7,-0.82)  {how do I structure\\ real work?};
    \node[q] at (8.55,-0.82) {how do I run a\\ whole project?};
    \node[q] at (11.4,-0.82) {what goes wrong?\\ what is my job?};
    % --- "you are here" marker above the current act ---
    \ifnum#1=1 \node[font=\scriptsize\bfseries, text=RedTitle] at (0,0.95)
        {you are here\ $\blacktriangledown$};\fi
    \ifnum#1=2 \node[font=\scriptsize\bfseries, text=RedTitle] at (2.85,0.95)
        {you are here\ $\blacktriangledown$};\fi
    \ifnum#1=3 \node[font=\scriptsize\bfseries, text=RedTitle] at (5.7,0.95)
        {you are here\ $\blacktriangledown$};\fi
    \ifnum#1=4 \node[font=\scriptsize\bfseries, text=RedTitle] at (8.55,0.95)
        {you are here\ $\blacktriangledown$};\fi
    \ifnum#1=5 \node[font=\scriptsize\bfseries, text=RedTitle] at (11.4,0.95)
        {you are here\ $\blacktriangledown$};\fi
    \ifnum#1=0 \node[font=\scriptsize\bfseries, text=green!45!black] at (5.7,0.95)
        {all five acts complete};\fi
    % --- bottom band: the three questions of the lecture ---
    \fill[blue!10]   (-1.3,-1.55) rectangle (4.15,-1.22);
    \fill[green!10]  (4.4,-1.55)  rectangle (9.85,-1.22);
    \fill[orange!15] (10.1,-1.55)  rectangle (12.7,-1.22);
    \node[barlab] at (1.425,-1.385)  {HOW IT WORKS};
    \node[barlab] at (7.125,-1.385)  {HOW TO USE IT WELL};
    \node[barlab] at (11.4,-1.385)   {YOUR ROLE};
    % --- dashed satellites: the two decks outside these five acts ---
    \node[sat] at (2.0,-2.2)
        {\textit{after 9.1:} optional interlude \textbf{9.1b} --- the agent landscape (MCP, A2A, benchmarks, safety)};
    \node[sat] at (9.8,-2.2)
        {\textit{after 9.5:} \textbf{9.6} wrap-up $\to$ \textbf{Lec10} full case studies (same morning)};
\end{tikzpicture}
\end{center}
}


\title[AI for Econ Research]{AI for Economic Research: Dynamic Models, Language, and Agents\\
Lecture 9.4: Project Organization --- The Researcher as Project Manager}

\begin{document}

\makebeamertitle

%=============================================================================
\begin{frame}{The Map: Act 4 --- The Management}
\textbf{This deck answers the fourth question: how do I run a whole project with this --- not just one task?}

\LecNineMap{4}

\vspace{0.1cm}
\footnotesize\textit{Route: the paradigm shift (researcher as project manager) $\to$ how HA-Ramsey was actually managed $\to$ the manager's instruments on disk (CLAUDE.md, skills, agents, rules, hooks) $\to$ your sub-agent RA team and cross-verification $\to$ remote/HPC execution.}
\end{frame}

%=============================================================================
\section{The Paradigm Shift}
%===========================================================

\sectiondivider{The Paradigm Shift}{From solo coder to project manager}

\begin{frame}{The Paradigm Shift: You Are the Project Manager}
\textbf{When execution becomes cheap, the scarce input is management: deciding what gets built, to what standard, by whom.}

\vspace{0.15cm}

\begin{center}
\begin{tikzpicture}[
    pm/.style={rectangle, rounded corners, draw=RedTitle!70, fill=blue!5,
        minimum width=1.92cm, minimum height=1.0cm, text width=1.84cm,
        align=center, font=\scriptsize},
    arr/.style={-{Stealth[length=2mm]}, thick, gray!60}
]
    \node[pm] (g)  at (0,0)     {\textbf{Goal}\\ \tiny one sentence + success criteria};
    \node[pm] (fp) at (2.17,0)  {\textbf{Final product}\\ \tiny paper, dataset, validated solver};
    \node[pm] (ig) at (4.34,0)  {\textbf{Intermediate goods}\\ \tiny spec, pseudocode, V0, batches};
    \node[pm] (qg) at (6.51,0)  {\textbf{Quality gate per good}\\ \tiny criterion set in advance};
    \node[pm] (as) at (8.68,0)  {\textbf{Assignment}\\ \tiny which session, agent, sub-agent};
    \node[pm] (cv) at (10.85,0) {\textbf{Cross-verifi\-cation}\\ \tiny producer never sole reviewer};
    \draw[arr] (g) -- (fp);
    \draw[arr] (fp) -- (ig);
    \draw[arr] (ig) -- (qg);
    \draw[arr] (qg) -- (as);
    \draw[arr] (as) -- (cv);
\end{tikzpicture}
\end{center}

\vspace{0.1cm}

\begin{itemize}
    \item Decompose the goal into \textbf{intermediate goods} --- real artifacts (spec, validated V0, batch CSV, review report), not vague ``progress''
    \item Decide the \textbf{quality gate} for each good \emph{before} it is produced (T3's acceptance criteria, now for everything, not just code)
    \item \textbf{Assign} each good to the right producer --- this session, a sub-agent, a different model tier, or yourself; then \textbf{cross-verify}: whoever produced it, someone else checks it
\end{itemize}

\vspace{0.05cm}
\textbf{T3 disciplined one artifact; T4 runs the factory.} Every object in this deck is one of the manager's instruments.
\end{frame}

\begin{frame}{A Managed Project in One Slide: HA-Ramsey's Milestones}
\textbf{The HA-Ramsey program is the PM chain instantiated: seven milestones, each a validated, preserved intermediate good.}

\vspace{0.15cm}

\begin{center}
\begin{tikzpicture}[
    m/.style={rectangle, rounded corners, draw=black!60, fill=VeryLightGray,
        minimum width=1.9cm, minimum height=1.05cm, text width=1.8cm,
        align=center, font=\tiny},
    md/.style={fill=green!10, draw=green!45!black},
    gate/.style={font=\tiny, text=black!60, align=center},
    arr/.style={-{Stealth[length=1.8mm]}, thick, gray!60}
]
    \node[m] (m1) at (0,0)    {\textbf{M1}\\refine the seed\\ \emph{zero new economics}};
    \node[m] (m2) at (2.2,0)  {\textbf{M2}\\new model in TeX\\ \emph{before any code}};
    \node[m] (m3) at (4.4,0)  {\textbf{M3}\\v1 prototype\\ \emph{proof of concept}};
    \node[m] (m4) at (6.6,0)  {\textbf{M4}\\full model v2\\ \emph{debugging $\to$ architecture}};
    \node[m, md] (m5) at (8.8,0) {\textbf{M5--M7}\\three new methods\\ \emph{discoveries}};
    \node[m] (fin) at (11.0,0){\textbf{Final}\\full paper +\\methods note};
    \draw[arr] (m1) -- (m2);
    \draw[arr] (m2) -- (m3);
    \draw[arr] (m3) -- (m4);
    \draw[arr] (m4) -- (m5);
    \draw[arr] (m5) -- (fin);
    \node[gate] at (1.1,-0.95)  {gate: TeX $\equiv$ code};
    \node[gate] at (3.3,-0.95)  {gate: economist audit};
    \node[gate] at (5.5,-0.95)  {gate: reproduces hand-derived checks};
    \node[gate] at (7.7,-0.95)  {gate: each v1 failure answered};
    \node[gate] at (9.9,-0.95)  {gate: written up, cited};
\end{tikzpicture}
\end{center}

\vspace{0.1cm}

\begin{itemize}
    \item Every milestone preserved as a citable artifact in the project's \texttt{Case\_Study/} package --- \emph{intermediate goods, literally}; no milestone skipped, no gate waived
    \item The green box is the payoff of management: the discoveries (Story 1 in T1) \emph{surfaced from gated iteration} --- they were not the plan
\end{itemize}

{\scriptsize \textit{Economics and figures: Lec10 T8. Source: HA-Ramsey case-study package.}}
\end{frame}

\begin{frame}{Write the Contract First: Who Owns What}
\textbf{Before any code, the HA-Ramsey project wrote down the division of labor --- T2's table, now with signatures.}

\vspace{0.2cm}

\begin{center}
\small
\begin{tabular}{|p{5.4cm}|p{5.8cm}|}
\hline
\textbf{AI owns (throughput)} & \textbf{Economist owns (judgment)} \\ \hline
Drafting and revising the LaTeX model note & Economic correctness of every equation \\ \hline
Translating math into vectorized PyTorch & Which simplifications are acceptable \\ \hline
Penalty/boundary implementation patterns & Catching conceptual bugs in timing and probabilities \\ \hline
Running diagnostics and summarizing sweeps & Validation benchmarks and when a result is credible \\ \hline
\end{tabular}
\end{center}

\vspace{0.15cm}

\begin{takeawaybox}
\footnotesize
\textbf{``AI supplies throughput; the economist supplies judgment''} --- stated \emph{before} any code was written. The contract is an assignment rule: which goods the AI may produce alone, and which gates only you may sign.
\end{takeawaybox}
\end{frame}

\begin{frame}{Quality Gates Catch What No Compiler Catches}
\textbf{Two conceptual bugs the economist caught in AI-drafted equations --- the code ran, the losses fell, the results would have been wrong. No compiler catches these:}

\begin{enumerate}
    \item \textbf{Wrong timing:} asset transitions written with next-period consumption --- but the budget constraint is a \emph{current-period} object
    \item \textbf{Wrong direction:} transition probabilities indexed by state occupied, not by \emph{direction of movement}
\end{enumerate}

\textbf{When v1 then failed numerically, each failure became v2 architecture:}
\begin{center}
\footnotesize
\begin{tabular}{|p{4.6cm}|p{6.2cm}|}
\hline
\textbf{v1 failure} & \textbf{v2 structural answer} \\ \hline
Naive sampling of a 5D state space & Adaptive sampling loops \\ \hline
Fragile feasible-set boundary learner & Normalized coordinates + robust geometry \\ \hline
Simulated paths escape the valid region & Explicit projection back + penalty term \\ \hline
\end{tabular}
\end{center}

\vspace{0.05cm}

\begin{cautionbox}
\footnotesize
\textbf{``Debugging becomes architecture'' only if someone with judgment reads the failures} --- the gate is where the economist earns authorship.
\end{cautionbox}
\end{frame}

%=============================================================================
\section{Project Organization}
%===========================================================

\sectiondivider{Project Organization}{The manager's instruments: files on disk}

\begin{frame}{The PM's Instruments Live on Disk}
\textbf{A manager who re-explains the standards every morning is not managing. Files on disk are how you delegate without re-explaining.}

\vspace{0.25cm}

\begin{itemize}
    \item Conventions live in files, not in your memory --- or the agent's
    \item Validation steps become reproducible lab procedure
    \item New sessions start from the project's actual state (T2: initialization reloads it all)
    \item Coauthors and students can inspect what the workflow is supposed to be
\end{itemize}

\vspace{0.25cm}

\textbf{Economist translation:} the harness turns tacit RA practice into written lab protocol --- and the protocol, not the chat, is what your name ultimately stands behind.
\end{frame}

\begin{frame}{Six Instruments, Six Pain Points}
\begin{center}
\small
\begin{tabular}{|p{2.1cm}|p{4.2cm}|p{4.1cm}|}
\hline
\textbf{Instrument} & \textbf{Management pain point it solves} & \textbf{Running example} \\ \hline
\texttt{CLAUDE.md} & re-explaining the project every session & keep the MEPS weights rule or Aiyagari benchmark visible \\ \hline
Skill & a repeated workflow improvised differently every time & paper-review pipeline with the same quality gates \\ \hline
Agent & review roles need specialization and parallelism & code reviewer vs.\ domain reviewer \\ \hline
Rule & some instructions must \emph{always} be on & every filled field carries a source URL \\ \hline
Verifier & output should pass named checks before it counts & rerun script, inspect diagnostics, confirm files exist \\ \hline
Provenance log & source discipline must survive beyond the chat & fellows dataset with field-level URLs and notes \\ \hline
\end{tabular}
\end{center}

\vspace{0.05cm}
\textbf{Reason for this slide:} these objects are not product trivia. Each one retires a coordination or credibility failure you would otherwise absorb personally.
\end{frame}

\begin{frame}[fragile]{Minimal Project Tree}
\begin{lstlisting}[style=pythonstyle,language={},basicstyle=\ttfamily\tiny]
my-research-project/
+-- CLAUDE.md
+-- .claude/
|   +-- skills/
|   |   +-- paper-review/
|   |       +-- SKILL.md
|   |       +-- references/
|   |           +-- referee-template.md
|   +-- agents/
|   |   +-- code-reviewer.md
|   |   +-- domain-reviewer.md
|   |   +-- verifier.md
|   +-- rules/
|   |   +-- iterative-workflow.md
|   +-- settings.json
+-- scripts/
+-- output/
+-- sources/
+-- quality_reports/
\end{lstlisting}

\textbf{This is enough to explain the workflow without burying students in infrastructure.} The exact research question can change; the organizational logic stays.
\end{frame}

\begin{frame}[fragile]{CLAUDE.md: The Standing Brief}
\texttt{CLAUDE.md} is the main shared instruction file for the project --- the brief a good manager writes once instead of repeating daily.

\vspace{0.2cm}

\begin{lstlisting}[style=pythonstyle,language={},basicstyle=\ttfamily\tiny,numbers=none]
# CLAUDE.md
Project: MEPS uninsurance pilot for lecture 9
Language: Python
Goal: harmonize 2012--2014 files and produce one validated figure
Conventions: no hardcoded paths; save outputs to output/
Checks: use survey weights; save the year map; report what passed
\end{lstlisting}

\vspace{0.2cm}

\textbf{Why students should care:} it captures the project context you would otherwise re-explain every session.

\vspace{0.1cm}
\footnotesize\textit{The fellows project runs the same play with public-web material: the seed CSV, the batch file, and the schema rules are ``the public-web analog of \texttt{CLAUDE.md}'' (Lec10 T4) --- project files holding the rules and state across sessions.}
\end{frame}

\begin{frame}{Three Memories, Three Writers}
\textbf{Three different persistence ideas should not be confused.}

\vspace{0.2cm}

\begin{center}
\small
\begin{tabular}{|p{2.4cm}|p{2.4cm}|p{5.3cm}|}
\hline
\textbf{Mechanism} & \textbf{Who writes it} & \textbf{Why it exists} \\ \hline
\texttt{CLAUDE.md} & you / team & durable project instructions and conventions \\ \hline
Auto memory & the tool & per-project learned patterns and corrections \\ \hline
Optional course log & you / course & explicit progress notes such as \texttt{progress.md} \\ \hline
\end{tabular}
\end{center}

\vspace{0.25cm}

\textbf{Course recommendation:} rely on \texttt{CLAUDE.md} for shared rules, treat auto memory as helpful but not authoritative, and keep explicit logs when you want students to inspect the evolution of a project.
\end{frame}

\begin{frame}{Skills, Agents, and Rules: Three Different Jobs}
\begin{center}
\begin{tikzpicture}[
    head/.style={rectangle, draw=RedTitle!70, fill=blue!8, rounded corners,
                 minimum width=3.4cm, minimum height=0.7cm,
                 text width=3.2cm, align=center, font=\small\bfseries},
    body/.style={rectangle, draw=black, fill=VeryLightGray, rounded corners,
                 minimum width=3.4cm, minimum height=1.55cm,
                 text width=3.2cm, align=center, font=\scriptsize, anchor=north},
    arr/.style={-{Stealth[length=1.8mm]}, thick, gray!70}
]
    \node[head] (hS) at (-4.0, 0) {Skills};
    \node[head] (hA) at ( 0.0, 0) {Agents};
    \node[head] (hR) at ( 4.0, 0) {Rules};

    \node[body] (bS) at (hS.south) {Callable workflows.\\
        Invoked by name, parameterized.\\
        \emph{Loaded on demand.}};
    \node[body] (bA) at (hA.south) {Specialist roles.\\
        Separate context, focused mandate.\\
        \emph{Spawned by the orchestrator.}};
    \node[body] (bR) at (hR.south) {Always-on guardrails.\\
        Applied automatically every session.\\
        \emph{Loaded at session start.}};

    \node[font=\tiny, gray!70] at (-2.0, -2.45) {load on demand $\rightarrow$};
    \node[font=\tiny, gray!70] at ( 2.0, -2.45) {$\leftarrow$ spawn on need};
    \node[font=\tiny, gray!70] at ( 0.0, -2.85) {Rules load \emph{first} every session; Skills and Agents enter the loop later.};
\end{tikzpicture}
\end{center}

\vspace{0.15cm}
\footnotesize\textbf{Loading order matters.} Rules are read before the first user message; Skills load when invoked; Agents launch with the rules and any inherited context as their system prompt.
\end{frame}

\begin{frame}{Plain English vs.\ Encoded Skill}
\begin{columns}[T]
\begin{column}{0.48\textwidth}
\textbf{Plain English}

\vspace{0.15cm}
\textit{Review this paper and tell me the top problems with exact quotes.}

\vspace{0.2cm}
\textbf{Use when}
\begin{itemize}
    \item Exploring
    \item Doing a one-off task
    \item Still deciding the workflow
\end{itemize}
\end{column}
\begin{column}{0.48\textwidth}
\textbf{Encoded skill}

\vspace{0.15cm}
\texttt{/paper-review draft\_A.pdf}

\vspace{0.2cm}
\textbf{Use when}
\begin{itemize}
    \item The workflow should repeat exactly
    \item A student team needs one standard procedure
    \item Validation must be consistent across versions
\end{itemize}
\end{column}
\end{columns}

\vspace{0.25cm}

\textbf{Exploration starts in English. Reproducibility lives in skills.}

\vspace{0.1cm}
\footnotesize\textit{Capstone note: Track D2 asks you to deliver exactly this --- a validated skill for a recurring research task.}
\end{frame}

\begin{frame}[fragile]{Skill File Structure}
\begin{lstlisting}[style=pythonstyle,language={},basicstyle=\ttfamily\tiny,numbers=none]
.claude/skills/paper-review/
  SKILL.md
  references/
    referee-template.md
\end{lstlisting}

\vspace{0.2cm}

\textbf{\texttt{SKILL.md}} stores the repeatable protocol:
\begin{itemize}
    \item what the command is called
    \item what steps it should follow
    \item what checks must pass before it reports success
\end{itemize}

\vspace{0.2cm}

\textbf{Reference files} store task-specific detail that would clutter the main instruction file.
\end{frame}

\begin{frame}{Skills Are Saved Research Playbooks}
\textbf{A skill is the right home for long, repeated procedures.} Keep always-loaded project context short; put domain workflows in reusable skills.

\vspace{0.1cm}
\begin{center}
\footnotesize
\begin{tabular}{|p{2.8cm}|p{4.0cm}|p{4.8cm}|}
\hline
\textbf{Skill} & \textbf{Reusable workflow} & \textbf{Validation hook} \\ \hline
\texttt{data-analysis} & Load, explore, regress, produce tables/figures & Output schema + script rerun \\ \hline
\texttt{review-paper} & Read paper, audit argument, generate referee objections & Structured review report \\ \hline
\texttt{weighted-survey} & Compute survey-weighted rates/means & Formula check against benchmark \\ \hline
\texttt{hpc-job} & Write \texttt{sbatch} script, submit, parse logs & Scheduler status + parser \\ \hline
\end{tabular}
\end{center}

\vspace{0.15cm}
\textbf{Portable idea:} Claude may call them \texttt{skills}; Codex may package them differently. The durable object is a saved playbook with triggers, allowed tools, steps, and acceptance tests.
\end{frame}

\begin{frame}[fragile]{Rules and Hooks: Encoding the Fellows Provenance Discipline}
\textbf{Guidance the agent should follow $\to$ a rule. A check that must fire mechanically $\to$ a hook. The fellows project needs both.}

\vspace{0.15cm}

\begin{columns}[T]
\begin{column}{0.52\textwidth}
\textbf{The rule} (\texttt{.claude/rules/provenance.md}):
\begin{lstlisting}[style=pythonstyle,language={},basicstyle=\ttfamily\tiny,numbers=none]
- Every filled field carries a source URL
  from an official or institutional page.
- Never infer gender, birth year, or living
  status from names, photos, or country.
- Record missing as missing; missingness
  is a result, not a failure.
\end{lstlisting}
\footnotesize Always-on: loaded at session start, batch after batch.
\end{column}
\begin{column}{0.45\textwidth}
\textbf{The hook} (\texttt{settings.json}):
\begin{lstlisting}[style=pythonstyle,language={},basicstyle=\ttfamily\tiny,numbers=none]
"PostToolUse": [{
  "matcher": "Write(data/*.csv)",
  "hooks": [{"type": "command",
   "command":
   "python checks/require_source_url.py"}]
}]
\end{lstlisting}
\footnotesize Mechanical: any CSV row missing \texttt{source\_url} is rejected --- no matter what the conversation said.
\end{column}
\end{columns}

\vspace{0.2cm}

\begin{center}
\footnotesize
\begin{tabular}{|p{1.6cm}|p{2.5cm}|p{6.0cm}|}
\hline
\textbf{Object} & \textbf{Fires} & \textbf{Use for} \\ \hline
Rules & during work & natural-language conventions and guardrails \\ \hline
Hooks & on tool events & mechanical enforcement: format, test, block, log \\ \hline
Settings & at session start & permissions and project-level behavior \\ \hline
\end{tabular}
\end{center}
\end{frame}

\begin{frame}[fragile]{Permissions and Hooks in \texttt{settings.json}}
\textbf{Two jobs:} \emph{permissions} define the blast radius; \emph{hooks} run mechanical checks at tool boundaries.

\vspace{0.15cm}
\begin{lstlisting}[style=pythonstyle,language={},basicstyle=\ttfamily\tiny,numbers=none]
{
  "permissions": {
    "allow": ["Bash(python *)", "Bash(pytest *)", "Read(*)", "Edit(*)"],
    "deny":  ["Bash(rm -rf *)", "Bash(git push --force *)",
              "Read(./.env)", "Read(./secrets/**)"]
  },
  "hooks": {
    "PostToolUse": [{"matcher": "Write|Edit|MultiEdit",
                     "hooks": [{"type": "command",
                                "command": ".claude/hooks/format-python.sh"}]}],
    "Stop":        [{"hooks": [{"type": "command",
                                "command": "pytest tests/ -q || true"}]}]
  }
}
\end{lstlisting}

\vspace{0.1cm}
\textbf{Hook types worth knowing:} \texttt{PreToolUse} (block risky commands), \texttt{PostToolUse} (format / lint), \texttt{Stop} (full test suite), \texttt{SubagentStop} (verify report format), \texttt{PreCompact} (save progress).
\end{frame}

% MiniCourse refinement: harness engineering vocabulary (specs / harness / skills) from M4_T7
\begin{frame}{Three Lifetimes: Per-Event, Per-Session, Forever}
\textbf{The harness is the umbrella name for what makes a project survive across sessions.}

\vspace{0.2cm}

\begin{center}
\small
\begin{tabular}{|l|l|p{6.0cm}|}
\hline
\textbf{Layer} & \textbf{Lifetime} & \textbf{Purpose} \\ \hline
\textbf{Hooks} (\texttt{settings.json}) & Per-event & Mechanical enforcement: format after edits, run tests on stop, block dangerous tools \\ \hline
\textbf{Project memory / rules} & Per-session & Bootstrap context. Project instructions and path-scoped guardrails \\ \hline
\textbf{Auto memory + progress files} & Cross-session & Durable learnings and explicit handoff notes you can inspect and edit \\ \hline
\end{tabular}
\end{center}

\vspace{0.12cm}
\footnotesize\textbf{Why this matters.} Choosing the right layer is half the design: if the rule must fire every time, a hook is safer than a natural-language reminder; if it must survive a fresh session, it belongs in memory or rules, not a transient prompt.

\vspace{0.04cm}
\textbf{Rule of thumb:} natural-language guidance $\to$ rules and skills. Mechanical enforcement $\to$ hooks.
\end{frame}

%===========================================================
% MiniCourse refinement: sub-agents / remote / HPC vocabulary from M4_T8
\section{Your RA Team}
%===========================================================

\sectiondivider{Your RA Team}{Sub-agents, assignment, and cross-verification}

\begin{frame}{Sub-Agents: Your Parallel RA Team}
\textbf{When the orchestrator spawns a sub-agent, the harness launches a separate agent instance with its own context window.}

\vspace{0.15cm}

\begin{itemize}
    \item Multiple sub-agents run \textbf{in parallel} --- like Python's \texttt{multiprocessing}
    \item Each receives an agent file as its system prompt and returns a structured report
    \item The specialist roles are the RA team you met in the project tree:
    \begin{itemize}
        \item \texttt{code-reviewer.md}: file hygiene, implementation quality, reproducibility
        \item \texttt{domain-reviewer.md}: benchmark logic, identification, economic consistency
        \item \texttt{verifier.md}: rerun the script, inspect diagnostics, confirm outputs exist
    \end{itemize}
\end{itemize}

\vspace{0.15cm}
\textbf{Aiyagari example.} After writing \texttt{aiyagari\_v1.py}, the orchestrator spawns
\begin{center}
\texttt{code-reviewer} \quad \textbf{$\Vert$} \quad \texttt{domain-reviewer}
\end{center}
in parallel; both finish in roughly one sequential run's wall-clock time. \textbf{Review time is halved.} \textit{Why separate roles? One prompt asking about everything does a mediocre job at each --- one RA checks style, another checks economics, another reruns the table.}
\end{frame}

\begin{frame}{Assignment and Cross-Verification}
\textbf{The manager's two questions for every intermediate good: who produces it, and who checks it?}

\begin{columns}[T]
\begin{column}{0.50\textwidth}
\textbf{Good fit for sub-agents:}
\begin{itemize}
    \item Code review (parallel: code, domain, verification)
    \item Lookups and web research that return short reports
    \item Independent unit-test generation per module
\end{itemize}
\end{column}
\begin{column}{0.48\textwidth}
\textbf{Bad fit:}
\begin{itemize}
    \item Tightly coupled work (each step depends on the previous)
    \item Anything you must steer interactively mid-task
    \item Sub-agents spawning sub-agents (depth causes drift)
\end{itemize}
\end{column}
\end{columns}

\vspace{0.15cm}

\textbf{Cross-verification: the producer is never the only reviewer.} Blast-radius rules of thumb: plan mode before granting execution on novel pipelines; each sub-agent's checklist \emph{spec'd}, not free-form; checkpoint validation between phases (the homotopy gates from T3).

\vspace{0.1cm}
\footnotesize\textit{Foreshadow for T5: parallel reviewers are parallel \emph{executionally}, not epistemically --- they read the same project files you wrote.}
\end{frame}

\begin{frame}{Orchestration at Research Scale}
\textbf{The HA-Ramsey lesson list ends with a management skill: to benefit the most, you need four things.}

\vspace{0.15cm}

\begin{enumerate}
    \item \textbf{A good seed} --- T3's criteria: works, documented, modular, points somewhere
    \item \textbf{A clear vision} --- the final product and the milestone ladder to it
    \item \textbf{Good judgment} --- the quality gates only an economist can sign
    \item \textbf{Orchestration} --- run \emph{different sessions and agents on different layers}, then integrate
\end{enumerate}

\vspace{0.15cm}

\textbf{What orchestration looked like in practice:}
\begin{itemize}
    \item Parallel sessions for formalization (TeX), implementation (PyTorch), auditing (code-vs-doc), sweeps, and write-up --- all feeding one preserved package
    \item 32+ configuration sweeps run and summarized by agents; the economist read summaries and decided at the milestones --- the job is \emph{integration}: the PM chain on a real paper
\end{itemize}

\vspace{0.05cm}
{\scriptsize \textit{Source: HA-Ramsey case-study package (Lec10 T8).}}
\end{frame}

%===========================================================
\section{Remote and HPC}
%===========================================================

\sectiondivider{Remote and HPC}{When the laptop is no longer the right host}

\begin{frame}{Remote/HPC: Move the Agent to the Compute}
\textbf{For heavy workloads the agent should run where the data, modules, and GPUs already are.}

\vspace{0.15cm}

\begin{center}
\begin{tikzpicture}[
    local/.style={rectangle, draw=blue!60, fill=blue!8, rounded corners,
                  minimum width=3.2cm, minimum height=1.5cm,
                  text width=3.0cm, align=center, font=\footnotesize},
    remote/.style={rectangle, draw=orange!70, fill=orange!10, rounded corners,
                   minimum width=3.6cm, minimum height=1.5cm,
                   text width=3.4cm, align=center, font=\footnotesize},
    compute/.style={rectangle, draw=green!50!black, fill=green!8, rounded corners,
                    minimum width=3.0cm, minimum height=1.0cm,
                    text width=2.8cm, align=center, font=\scriptsize},
    arr/.style={-{Stealth[length=2mm]}, thick, gray!60}
]
    \node[local] (mac) {\textbf{Laptop (you)}\\editor / terminal front end};
    \node[remote, right=1.3cm of mac] (login) {\textbf{Cluster login node}\\agent runs here\\reads/edits code, submits jobs};
    \node[compute, below=0.7cm of login] (cn) {\textbf{Compute nodes}\\MPI / GPU / long simulations};

    \draw[arr] (mac) -- node[above,font=\tiny]{SSH / VS Code remote} (login);
    \draw[arr] (login) -- node[right,font=\tiny]{\texttt{sbatch}, \texttt{srun}} (cn);
    \draw[arr] (cn) -- ++(2.0,0) -- ++(0,1.6) -- (login.east);
    \node[font=\tiny, gray!70, right=0.05cm of cn, xshift=2.1cm, yshift=0.8cm] {logs, results};
\end{tikzpicture}
\end{center}

\vspace{0.15cm}

\begin{itemize}
    \item With a local agent, cluster files must be copied back and forth --- easy for toy code, fragile for research
    \item With a remote agent, project files stay in the remote worktree; your laptop is only the front end
\end{itemize}

\vspace{0.05cm}
\textbf{Rule:} if the data, modules, GPUs, and job logs live on the cluster, install the agent on the cluster.
\end{frame}

\begin{frame}{Login Node vs.\ Compute Node: The Hard Boundary}
\textbf{What the agent can do on a login node:}
\begin{itemize}
    \item Read/edit code, specs, context files, and job scripts
    \item Run tiny smoke tests: import checks, unit tests, \texttt{--dry-run}, small-grid V0
    \item Submit jobs, monitor logs, parse errors, and revise scripts
\end{itemize}

\vspace{0.15cm}
\textbf{What must go through the scheduler:}
\begin{itemize}
    \item MPI jobs, GPU training, large simulations, long calibrations
    \item Anything that consumes many cores, much memory, or more than a few minutes
\end{itemize}

\bigskip
\begin{center}
\textbf{Research rule:} the agent may orchestrate HPC work; it should not turn the login node into a workstation.
\end{center}
\end{frame}

\begin{frame}{Bridge: Management Assumes Trust --- Verify It}
\textbf{You now delegate through files and a sub-agent team. Delegation works only if you can audit what was actually done --- and know when to distrust the polish.}

\vspace{0.25cm}

\begin{itemize}
    \item Act 4 gave you the management system: goods, gates, assignment, cross-verification, and the instruments that encode them
    \item The final question: \emph{``what can go wrong, and what is my job?''}
    \item Act 5 answers with three moves: \textbf{trace} what the agent did, catalog the \textbf{failure modes} the harness cannot prevent, and confront the \textbf{Context Dilemma} --- the same context that improves the answer also pulls it toward you
\end{itemize}

\vspace{0.3cm}
\footnotesize\textit{Arc: T1 Framing $\to$ T2 Under the Hood $\to$ T3 Homotopy $\to$ \textbf{T4 Project Org} $\to$ T5 Mistakes/Context $\to$ T6 Wrap.}
\end{frame}

\end{document}
